\documentclass[12pt]{article}
\usepackage[utf8]{inputenc}
\usepackage[spanish]{babel}
\usepackage{amsmath, mathrsfs}
\usepackage{amsfonts}
\usepackage{amssymb}
\usepackage{geometry}
\usepackage{fancyhdr}
\usepackage{parskip}
\usepackage{float}
\usepackage{hyperref}
\usepackage{tikz}
\usetikzlibrary{arrows.meta,positioning}
\AtBeginEnvironment{tikzpicture}{\shorthandoff{<>}}

\geometry{a4paper, margin=1in}

\pagestyle{fancy}
\fancyhf{}
\setlength{\headheight}{14.5pt}
\lhead{Astroestadística }
\chead{Problem Set 6}
\rhead{Problemas}


\begin{document}


\section*{Problema 1}
    
    \textbf{Generación de números aleatorios.} Utilice los métodos de inversión y/o de aceptación y rechazo para generar un conjunto de números aleatorios que sigan las siguientes distribuciones:
    \begin{itemize}
        \item Distribución de Weibull con $\alpha = 1.5$ y $\beta = 1$.
        \item Distribución exponencial.
        \item Distribución T con 10 grados de libertad.
        \item $f_X(x) = 3x^2, \; 0 < x < 1$.
    \end{itemize}
    
\section*{Problema 2}
    \textbf{Integración de Monte Carlo.} Utilice el método de Monte Carlo para calcular las integrales de las siguientes funciones (elija los límites de integración):
    \begin{itemize}
        \item $x^2$
        \item $e^x$
        \item $x^4 e^{-x}$
        \item $\frac{(e^{-x})}{(1 + (x - 1)^2)}$
    \end{itemize}

\section*{Problema 3}
    \textbf{Desviación de la normal.} Utilice el método de Monte Carlo para estudiar la desviación de los datos respecto a la normal. Tome $\mu = 0$ y $\sigma^2 = 1$ y genere una muestra de 100 números con una distribución $N(\mu, \sigma^2)$. De esta muestra, calcule $\bar{x}_1$ y $s_1^2$. Ahora, haga $\mu = \bar{x}_1$ y $\sigma^2 = s_1^2$; genere una muestra de 100 números y calcule ahora $\bar{x}_2$ y $s_2^2$. Continúe de esta manera, generando secuencias $\{\bar{x}_k\}$ y $\{s_k^2\}$. Describa este proceso estocástico.



\end{document}